Discrete choice models are central to transportation, marketing, operations, and service-system research, but researchers who need both econometric inference and modern tensor computation still face a software gap.
Mature packages such as Biogeme, Apollo, and \texttt{mlogit} provide rich model syntax and inference, whereas tensor frameworks provide vectorized automatic differentiation but do not expose a discrete-choice-oriented estimation and benchmarking workflow.
We introduce TorchDCM, a PyTorch-first package for discrete choice model estimation, inference, post-estimation analysis, and reproducible estimator benchmarking.
TorchDCM implements a model zoo that includes multinomial logit, nested logit, cross-nested logit, mixed logit, WTP-space mixed logit, ordered response models, latent class models, and hybrid choice components, while retaining familiar econometric outputs such as covariance matrices, standard errors, willingness-to-pay measures, elasticities, and probability predictions.
The package is paired with a public benchmark system that aligns data, model specifications, starting values, covariance definitions, probability calculations, and runtime accounting against Biogeme, Apollo, SciPy, \texttt{mlogit}, \texttt{gmnl}, and \texttt{xlogit}.
On current public benchmarks, TorchDCM matches reference estimators to numerical tolerance for full-estimation MNL, nested logit, cross-nested logit, ordered logit, ordered probit, and an aligned Swissmetro mixed-logit case against Biogeme, and it verifies additional mixed-logit likelihood kernels using shared simulation draws.
Across the full Swissmetro public-data benchmark after common filtering, TorchDCM completes full-estimation MNL in 0.030 seconds versus 1.188 seconds for the fastest econometric reference, nested logit in 0.069 seconds versus 2.100 seconds, and cross-nested logit in 1.108 seconds versus 3.870 seconds.
On an NHTS 2022 public travel-survey mode-choice benchmark, TorchDCM estimates a 27,375-trip, 16-parameter MNL in 0.121 seconds versus 1.735 seconds for Biogeme.
These results position TorchDCM as open infrastructure for choice-model research that needs differentiable computation, econometric reporting, and transparent estimator parity checks.
